\documentclass[aspectratio=169,10pt]{beamer}
\makeatletter\def\input@path{{../estilo/}{estilo/}}\makeatother
\usepackage{devmobile}
\modulo{13}{Desenvolvimento Multiplataforma}{Flutter, React Native e Kotlin Multiplatform: arquitetura, código comparado e critérios de escolha}

\begin{document}
\capa

\begin{frame}{Objetivos do módulo}
  \begin{conceito}[Ao final deste módulo você será capaz de]
    \begin{itemize}
      \item Explicar como Flutter, React Native e Kotlin Multiplatform chegam ao Android e ao iOS.
      \item Ler e escrever uma tela simples em Dart/Flutter e em TypeScript/React Native.
      \item Estruturar um projeto Kotlin Multiplatform com código comum e específico.
      \item Mapear os conceitos de Compose para os equivalentes em cada framework.
      \item Escolher a abordagem adequada a um projeto e a uma equipe.
    \end{itemize}
  \end{conceito}
  \begin{dica}[Você já sabe mais do que imagina]
    Os três frameworks usam \textbf{UI declarativa}, \textbf{estado} e \textbf{componentes}, os mesmos conceitos dos Módulos 04--06.
  \end{dica}
\end{frame}

\sumario

% =============================================================================
\section{Panorama}

\begin{frame}{Três estratégias de compartilhamento}
  \begin{center}
  \begin{tikzpicture}[node distance=2mm]
    % Flutter
    \node[caixaambarforte,minimum width=3.6cm] (f1) {Código Dart (UI + lógica)};
    \node[caixa,minimum width=3.6cm,below=of f1] (f2) {Motor próprio (Impeller)\\ \tiny desenha cada pixel};
    \node[caixa,minimum width=1.7cm,below=of f2.south west,anchor=north west,fill=cinzaclaro] (f3) {Android};
    \node[caixa,minimum width=1.7cm,below=of f2.south east,anchor=north east,fill=cinzaclaro] (f4) {iOS};
    \node[above=1mm of f1,font=\small\bfseries,text=petroleo] {Flutter};
    % RN
    \node[caixaambarforte,minimum width=3.6cm,right=10mm of f1] (r1) {Código JS/TS (UI + lógica)};
    \node[caixa,minimum width=3.6cm,below=of r1] (r2) {Ponte JSI\\ \tiny componentes nativos};
    \node[caixa,minimum width=1.7cm,below=of r2.south west,anchor=north west,fill=cinzaclaro] (r3) {Android};
    \node[caixa,minimum width=1.7cm,below=of r2.south east,anchor=north east,fill=cinzaclaro] (r4) {iOS};
    \node[above=1mm of r1,font=\small\bfseries,text=petroleo] {React Native};
    % KMP
    \node[caixaescura,minimum width=3.6cm,right=10mm of r1] (k1) {Kotlin comum (lógica)};
    \node[caixa,minimum width=3.6cm,below=of k1] (k2) {UI: Compose Multiplatform\\ \tiny ou nativa (SwiftUI)};
    \node[caixa,minimum width=1.7cm,below=of k2.south west,anchor=north west,fill=cinzaclaro] (k3) {Android};
    \node[caixa,minimum width=1.7cm,below=of k2.south east,anchor=north east,fill=cinzaclaro] (k4) {iOS};
    \node[above=1mm of k1,font=\small\bfseries,text=petroleo] {Kotlin Multiplatform};
  \end{tikzpicture}
  \end{center}
  \begin{columns}[T]
    \begin{column}{0.32\textwidth}
      \begin{conceito}[Flutter]
        Compartilha \textbf{tudo}; a interface é desenhada pelo próprio motor, idêntica nas plataformas.
      \end{conceito}
    \end{column}
    \begin{column}{0.32\textwidth}
      \begin{conceito}[React Native]
        Compartilha \textbf{tudo}; os componentes viram widgets nativos de cada sistema.
      \end{conceito}
    \end{column}
    \begin{column}{0.32\textwidth}
      \begin{conceito}[KMP]
        Compartilha \textbf{o que fizer sentido}: só a lógica ou também a UI.
      \end{conceito}
    \end{column}
  \end{columns}
\end{frame}

\begin{frame}{Comparação}
  \footnotesize
  \begin{tabularx}{\textwidth}{@{}l L L L@{}}
    \toprule
    & \cabecalho{Flutter} & \cabecalho{React Native} & \cabecalho{Kotlin Multiplatform} \\
    \midrule
    Linguagem & Dart & JavaScript / TypeScript & Kotlin \\
    Mantenedor & Google & Meta e comunidade & JetBrains (com apoio do Google) \\
    Interface & widgets próprios desenhados pelo motor & componentes nativos & Compose Multiplatform ou nativa \\
    Desenvolvimento & \emph{hot reload} & \emph{fast refresh} & \emph{hot reload} (Compose) \\
    Acesso nativo & \emph{plugins} e \emph{platform channels} & módulos nativos (Turbo Modules) & direto (\cod{expect}/\cod{actual}) \\
    Outras plataformas & web, desktop & web (parcial) & desktop, web, servidor \\
    Para quem sabe Kotlin & aprender Dart & aprender JS/TS e React & continua em Kotlin \\
    \bottomrule
  \end{tabularx}
  \vspace{4pt}
  \begin{dica}
    Todos são usados em produção por grandes empresas. Escolha pelo \textbf{contexto} (equipe, código existente, requisitos), não por modismo.
  \end{dica}
\end{frame}

% =============================================================================
\section{Flutter}

\begin{frame}[fragile]{Flutter: tudo é widget}
  \begin{columns}[T]
    \begin{column}{0.58\textwidth}
\begin{codigo}[title=lib/painel.dart]{Dart}
class Painel extends StatefulWidget {
  const Painel({super.key});
  @override
  State<Painel> createState() => _PainelState();
}

class _PainelState extends State<Painel> {
  double temp = 24.0;                 // estado

  @override
  Widget build(BuildContext context) => Center(
    child: Column(mainAxisSize: MainAxisSize.min,
      children: [
        Text('${temp.toStringAsFixed(1)} °C'),
        ElevatedButton(
          onPressed: () => setState(() => temp += 0.5),
          child: const Text('Simular'),
        ),
      ]),
  );
}
\end{codigo}
    \end{column}
    \begin{column}{0.4\textwidth}
      \begin{conceito}[Conceitos]
        \begin{itemize}
          \item \cod{StatelessWidget}: sem estado próprio.
          \item \cod{StatefulWidget} + \cod{State}: com estado; \cod{setState} pede um novo \cod{build}.
          \item Layout com widgets: \cod{Column}, \cod{Row}, \cod{Padding}, \cod{Center}.
        \end{itemize}
      \end{conceito}
      \begin{dica}[Compose $\leftrightarrow$ Flutter]
        \scriptsize
        \begin{tabular}{@{}l l@{}}
          \cod{@Composable} & \cod{Widget.build} \\
          \cod{mutableStateOf} & \cod{setState} \\
          \cod{Modifier.padding} & \cod{Padding(...)} \\
          \cod{LazyColumn} & \cod{ListView.builder} \\
          \cod{ViewModel} & \cod{ChangeNotifier}/Bloc \\
        \end{tabular}
      \end{dica}
    \end{column}
  \end{columns}
\end{frame}

\begin{frame}[fragile]{Flutter: ecossistema e estrutura}
  \begin{columns}[T]
    \begin{column}{0.46\textwidth}
\begin{terminal}[title=Começando]
flutter create labmonitor_flutter
cd labmonitor_flutter
flutter run          # r = hot reload
flutter pub add http riverpod
flutter build appbundle
\end{terminal}
\begin{codigo}[title=pubspec.yaml (trecho)]{bash}
dependencies:
  flutter:
    sdk: flutter
  http: ^1.2.0
  flutter_riverpod: ^2.5.0
\end{codigo}
      \begin{dica}
        Pacotes em \url{https://pub.dev}. Confira as versões atuais antes de copiar.
      \end{dica}
    \end{column}
    \begin{column}{0.52\textwidth}
      \scriptsize
      \renewcommand{\arraystretch}{1.1}
      \begin{tabular}{@{}l l@{}}
        \toprule
        \cabecalho{Necessidade} & \cabecalho{Pacotes comuns} \\
        \midrule
        Estado & Provider, Riverpod, Bloc \\
        Rede & \cod{http}, \cod{dio} \\
        Banco local & \cod{drift}, \cod{sqflite} \\
        Firebase & FlutterFire \\
        BLE & \cod{flutter\_blue\_plus} \\
        MQTT & \cod{mqtt\_client} \\
        Navegação & \cod{go\_router} \\
        \bottomrule
      \end{tabular}
      \begin{conceito}[Pontos fortes e limites]
        UI idêntica em todas as plataformas, \emph{hot reload} rápido e ótima documentação. O motor embutido deixa o APK maior que o nativo.
      \end{conceito}
    \end{column}
  \end{columns}
\end{frame}

% =============================================================================
\section{React Native}

\begin{frame}[fragile]{React Native: componentes e hooks}
  \begin{columns}[T]
    \begin{column}{0.56\textwidth}
\begin{codigo}[title=app/index.tsx (Expo)]{TypeScript}
import { useState } from 'react';
import { View, Text, Button, StyleSheet } from 'react-native';

export default function Painel() {
  const [temp, setTemp] = useState(24.0);   // estado

  return (
    <View style={styles.centro}>
      <Text style={styles.valor}>{temp.toFixed(1)} °C</Text>
      <Button title="Simular"
        onPress={() => setTemp(t => t + 0.5)} />
    </View>
  );
}

const styles = StyleSheet.create({
  centro: { flex: 1, alignItems: 'center',
            justifyContent: 'center' },
  valor: { fontSize: 48, fontWeight: 'bold' },
});
\end{codigo}
    \end{column}
    \begin{column}{0.42\textwidth}
      \begin{conceito}[Conceitos]
        \begin{itemize}
          \item Componentes são funções que devolvem JSX.
          \item \cod{useState}, \cod{useEffect}: \emph{hooks} (como \cod{remember} e \cod{LaunchedEffect}).
          \item Estilo parecido com CSS e layout \emph{flexbox}.
          \item Nova arquitetura (Fabric, Turbo Modules, JSI).
        \end{itemize}
      \end{conceito}
      \begin{dica}[Expo]
        Caminho recomendado para começar: \cod{npx create-expo-app}, rotas por arquivos e o app \emph{Expo Go}.
      \end{dica}
    \end{column}
  \end{columns}
\end{frame}

% =============================================================================
\section{Kotlin Multiplatform}

\begin{frame}[fragile]{Estrutura de um projeto KMP}
  \begin{columns}[T]
    \begin{column}{0.46\textwidth}
\begin{codigo}[title=Projeto KMP]{bash}
labmonitor-kmp/
|-- composeApp/
|   `-- src/
|       |-- commonMain/   UI Compose comum
|       |-- androidMain/  MainActivity
|       `-- iosMain/      MainViewController
|-- shared/
|   `-- src/
|       |-- commonMain/   modelos, repositórios,
|       |                 casos de uso, Ktor
|       |-- androidMain/  implementações Android
|       `-- iosMain/      implementações iOS
`-- iosApp/               projeto Xcode (Swift)
\end{codigo}
      \begin{dica}
        O \textbf{Kotlin Multiplatform Wizard} (\url{https://kmp.jetbrains.com}) gera essa estrutura pronta.
      \end{dica}
    \end{column}
    \begin{column}{0.52\textwidth}
      \centering
      \begin{tikzpicture}[node distance=3mm]
        \node[caixaescura,minimum width=5cm] (c) {\cod{commonMain}\\ \tiny Kotlin puro + bibliotecas KMP};
        \node[caixaambar,minimum width=2.4cm,below=6mm of c.south west,anchor=north west] (a) {\cod{androidMain}\\ \tiny JVM / Android SDK};
        \node[caixaambar,minimum width=2.4cm,below=6mm of c.south east,anchor=north east] (i) {\cod{iosMain}\\ \tiny Kotlin/Native};
        \draw[seta] (c) -- (a); \draw[seta] (c) -- (i);
        \node[caixa,below=of a,minimum width=2.4cm,fill=cinzaclaro] (ap) {\arq{.aab} Android};
        \node[caixa,below=of i,minimum width=2.4cm,fill=cinzaclaro] (ip) {\arq{.framework} iOS};
        \draw[seta] (a) -- (ap); \draw[seta] (i) -- (ip);
      \end{tikzpicture}
      \scriptsize
      \begin{tabularx}{\linewidth}{@{}l L@{}}
        \toprule
        \cabecalho{Camada} & \cabecalho{Bibliotecas KMP} \\
        \midrule
        Rede & Ktor Client + kotlinx.serialization \\
        Banco & Room (KMP), SQLDelight, DataStore \\
        Estado & ViewModel (lifecycle KMP), coroutines \\
        DI & Koin \\
        UI & Compose Multiplatform \\
        \bottomrule
      \end{tabularx}
    \end{column}
  \end{columns}
\end{frame}

\begin{frame}[fragile]{\cod{expect} e \cod{actual}}
  \begin{columns}[T]
    \begin{column}{0.54\textwidth}
\begin{codigo}[title=commonMain/Plataforma.kt]{Kotlin}
expect fun nomePlataforma(): String
\end{codigo}
\begin{codigo}[title=androidMain/Plataforma.android.kt]{Kotlin}
actual fun nomePlataforma() =
  "Android ${Build.VERSION.SDK_INT}"
\end{codigo}
\begin{codigo}[title=iosMain/Plataforma.ios.kt]{Kotlin}
import platform.UIKit.UIDevice

actual fun nomePlataforma() =
  UIDevice.currentDevice.systemName() + " " +
  UIDevice.currentDevice.systemVersion
\end{codigo}
\begin{codigo}[title=commonMain: uso]{Kotlin}
Text("Rodando em ${nomePlataforma()}")
\end{codigo}
    \end{column}
    \begin{column}{0.44\textwidth}
      \begin{conceito}[O mecanismo]
        \cod{expect} declara \textbf{o quê} no código comum; cada plataforma fornece o \cod{actual}, \textbf{como}, usando suas APIs nativas (Android SDK ou frameworks da Apple).
      \end{conceito}
      \begin{dica}[Prefira interfaces]
        Para dependências maiores (BLE, sensores), declare uma \cod{interface} no código comum e injete implementações por plataforma com Koin. Fica mais fácil de testar.
      \end{dica}
      \begin{atencao}
        Compilar e rodar a parte iOS exige \textbf{macOS com Xcode}.
      \end{atencao}
    \end{column}
  \end{columns}
\end{frame}

\begin{frame}[fragile]{Compartilhando a lógica do LabMonitor}
  \begin{columns}[T]
    \begin{column}{0.56\textwidth}
\begin{codigo}[title=shared/commonMain/ClimaRepository.kt]{Kotlin}
class ClimaRepository(private val http: HttpClient) {
  suspend fun atual(lat: Double, lon: Double): ClimaDto =
    http.get("https://api.open-meteo.com/v1/forecast") {
      parameter("latitude", lat)
      parameter("longitude", lon)
      parameter("current", "temperature_2m")
    }.body()
}

class ClimaViewModel(
  private val repo: ClimaRepository
) : ViewModel() {                 // lifecycle KMP
  val estado = MutableStateFlow<Double?>(null)
  fun carregar() = viewModelScope.launch {
    estado.value = repo.atual(-2.53, -44.30)
      .current.temperatura
  }
}
\end{codigo}
    \end{column}
    \begin{column}{0.42\textwidth}
      \begin{conceito}[O que mudou em relação ao Android]
        Quase nada: Retrofit $\rightarrow$ \textbf{Ktor}, Hilt $\rightarrow$ \textbf{Koin}. DTOs, \cod{Flow}, \cod{ViewModel}, casos de uso e testes continuam iguais.
      \end{conceito}
      \begin{dica}[Migração gradual]
        Um app Android existente pode mover \textbf{primeiro} os modelos e casos de uso para um módulo \cod{shared}, sem reescrever a UI. É a forma mais segura de adotar KMP.
      \end{dica}
    \end{column}
  \end{columns}
\end{frame}

% =============================================================================
\section{Como escolher}

\begin{frame}{A mesma tela nos três mundos}
  \footnotesize
  \begin{tabularx}{\textwidth}{@{}l L L L@{}}
    \toprule
    \cabecalho{Conceito} & \cabecalho{Compose (Kotlin)} & \cabecalho{Flutter (Dart)} & \cabecalho{React Native (TS)} \\
    \midrule
    Unidade de UI & função \cod{@Composable} & classe \cod{Widget} & função componente \\
    Estado local & \cod{remember \{ mutableStateOf() \}} & \cod{State} + \cod{setState} & \cod{useState} \\
    Efeito ao entrar & \cod{LaunchedEffect} & \cod{initState} & \cod{useEffect} \\
    Lista longa & \cod{LazyColumn} & \cod{ListView.builder} & \cod{FlatList} \\
    Layout & \cod{Column}/\cod{Row}/\cod{Box} & \cod{Column}/\cod{Row}/\cod{Stack} & \cod{View} + flexbox \\
    Estado da tela & \cod{ViewModel} + \cod{StateFlow} & Riverpod / Bloc & Context, Zustand, Redux \\
    Assíncrono & coroutines / \cod{Flow} & \cod{Future} / \cod{Stream} & \cod{Promise} / \cod{async} \\
    \bottomrule
  \end{tabularx}
  \vspace{4pt}
  \begin{destaque}
    Os conceitos se transferem: quem domina um framework declarativo aprende o outro em semanas.
  \end{destaque}
\end{frame}

\begin{frame}{Critérios de decisão}
  \begin{columns}[T]
    \begin{column}{0.5\textwidth}
      \begin{conceito}[Perguntas a fazer]
        \begin{enumerate}
          \item Precisa de iOS? Quando?
          \item Que linguagem a equipe já domina?
          \item Há código Android/Kotlin existente?
          \item O app depende muito de hardware (BLE, câmera, sensores)?
          \item A UI deve ser idêntica ou seguir o padrão de cada sistema?
          \item Há também web ou desktop no horizonte?
        \end{enumerate}
      \end{conceito}
    \end{column}
    \begin{column}{0.48\textwidth}
      \scriptsize
      \begin{tabularx}{\linewidth}{@{}L l@{}}
        \toprule
        \cabecalho{Cenário} & \cabecalho{Sugestão} \\
        \midrule
        Só Android, hardware intenso & Kotlin nativo \\
        App Android existente indo para iOS & KMP \\
        Equipe web/React & React Native \\
        UI de marca idêntica em tudo, equipe nova & Flutter \\
        Protótipo rápido para as duas lojas & Flutter ou RN + Expo \\
        \bottomrule
      \end{tabularx}
      \begin{dica}[Para o LabMonitor]
        Como o app usa BLE e já está em Kotlin, o caminho natural para o iOS é \textbf{KMP}: compartilhar modelos, repositórios e casos de uso, e escolher entre Compose Multiplatform e SwiftUI para a UI.
      \end{dica}
    \end{column}
  \end{columns}
\end{frame}

% =============================================================================
\section{Fechamento}

\begin{frame}{Resumo}
  \begin{resumo}
    \begin{itemize}
      \item Flutter desenha a própria UI com Dart; React Native usa componentes nativos com JS/TS; KMP compartilha Kotlin e deixa a UI a critério.
      \item Os três são declarativos: estado, componentes e efeitos se mapeiam diretamente para Compose.
      \item KMP usa \cod{commonMain} + \cod{expect}/\cod{actual} e bibliotecas como Ktor, Room, Koin e Compose Multiplatform.
      \item iOS sempre exige macOS e Xcode para compilar e publicar.
      \item A escolha depende da equipe, do código existente, do hardware e das plataformas-alvo.
    \end{itemize}
  \end{resumo}
\end{frame}

\begin{frame}{Exercícios}
  \begin{exercicio}[Exercício 1 --- Flutter]
    Reescreva em Flutter a tela de temperatura da Open-Meteo (Módulo 08) com \cod{http} e um \cod{FutureBuilder} mostrando carregando, sucesso e erro.
  \end{exercicio}
  \begin{exercicio}[Exercício 2 --- React Native]
    Crie com Expo uma lista de equipamentos com \cod{FlatList} e uma tela de detalhe usando o roteamento por arquivos.
  \end{exercicio}
  \begin{exercicio}[Exercício 3 --- Comparação]
    Meça, para o app do exercício 1 e para o equivalente em Compose, o tamanho do APK release e o tempo de abertura. Discuta os resultados.
  \end{exercicio}
  \begin{projeto}[LabMonitor --- Etapa 13]
    Extraia modelos, DTOs e o \cod{VerificarAlertasUseCase} para um módulo \cod{shared} KMP com testes em \cod{commonTest}. Opcional: tela de painel em Compose Multiplatform rodando no desktop.
  \end{projeto}
\end{frame}

\begin{frame}{Para aprofundar}
  \small
  \begin{itemize}
    \item Flutter: \url{https://docs.flutter.dev} \quad Flutter para devs Android: \url{https://docs.flutter.dev/get-started/flutter-for/android-devs}
    \item React Native: \url{https://reactnative.dev} \quad Expo: \url{https://docs.expo.dev}
    \item Kotlin Multiplatform: \url{https://kotlinlang.org/docs/multiplatform.html}
    \item KMP no Android: \url{https://developer.android.com/kotlin/multiplatform}
    \item Compose Multiplatform: \url{https://www.jetbrains.com/compose-multiplatform/}
    \item Ktor Client: \url{https://ktor.io/docs/client-create-new-application.html}
  \end{itemize}
\end{frame}

\encerramento{Módulo 14 --- Tópicos avançados e projeto final: IA no dispositivo, segurança, monetização e carreira.}

\end{document}
